% Volume I, appendice «Elementi di statica grafica»: figura fig:tre_forze_sgr.
% Richiamata con \FiguraSGR{figura-sgr-tre-forze}; didascalia e label nell'appendice.
% Generata con gen_sgr.py (coordinate numeriche; stili di Preambolo/tikz-pol e tikz-sgr).
%
% Reazioni dell'esercizio ex:anello_vincolato_TA (capitolo Sistemi a connessioni eterogenee), in unita' l:
% A(0,0), B(2,5; 0), C(0,5; 1), D(1,5; 1), E(0,5; 2), G(1; 1); bielle AC e CE, travi CD ed EDB, glifo in D,
% cerniere a terra in A e B. Carico p sulla proiezione orizzontale di ED (risultante p l su x = l) e coppia
% oraria p l^2 in G: forza e coppia equivalgono alla forza p l sulla verticale x = 2 l.
% La reazione in A ha la direzione della biella AC; le tre rette concorrono in K(2; 4), per cui passa la
% reazione in B. Triangolo delle forze: r_A = (0.100,0.200) p l, r_B = (-0.100,0.800) p l (verificato: V_B = 4/5, V_A = 1/5,
% H_A = 1/10, H_B = 1/10 verso sinistra, in p l).
% Struttura in scala 1.35 cm per l; triangolo delle forze in scala 3.6 cm per p l.
% (a) struttura, forza equivalente, rette d'azione delle reazioni e punto K; (b) triangolo delle forze.
\begin{tikzpicture}[font=\footnotesize]
\begin{scope}[x=1.35cm,y=1.35cm]
\draw[polretta] (2.000,-0.350) -- (2.000,4.450) node[anchor=south,inner sep=1.5pt,black] {$a$};
\draw[polcostruzione] (0.500,1.000) -- (2.089,4.179);
\draw[polcostruzione] (2.500,0.000) -- (1.963,4.298);
\node[anchor=east,inner sep=3pt] at (0.250,0.500) {$a_A$};
\node[anchor=west,inner sep=3pt] at (2.275,1.800) {$a_B$};
\draw[sgrbiella] (0.000,0.000) -- (0.500,1.000) -- (0.500,2.000);
\draw[sgrtrave] (0.500,2.000) -- (2.500,0.000);
\draw[sgrtrave] (0.500,1.000) -- (1.34,1.0);
\draw[line width=1.3pt] (1.277,1.060) -- (1.560,0.777);
\draw[fill=white,line width=.5pt] (1.388,1.030) circle (0.042);
\draw[fill=white,line width=.5pt] (1.530,0.888) circle (0.042);
\pic at (0.000,0.000) {polcerniera};
\pic at (2.500,0.000) {polcerniera};
\pic at (0.500,1.000) {sgrnodo};
\pic at (0.500,2.000) {sgrnodo};
\node[anchor=east,inner sep=4pt] at (0.000,0.000) {$A$};
\node[anchor=west,inner sep=4pt] at (2.500,0.000) {$B$};
\node[anchor=east,inner sep=4pt] at (0.500,1.000) {$C$};
\node[anchor=east,inner sep=4pt] at (0.500,2.000) {$E$};
\node[anchor=south west,inner sep=3pt] at (1.500,1.000) {$D$};
\draw[line width=.4pt] (0.5,2.28) -- (1.5,2.28) (0.5,2.55) -- (1.5,2.55);
\draw[polcarico] (0.500,2.55) -- (0.500,2.28);
\draw[polcarico] (0.667,2.55) -- (0.667,2.28);
\draw[polcarico] (0.833,2.55) -- (0.833,2.28);
\draw[polcarico] (1.000,2.55) -- (1.000,2.28);
\draw[polcarico] (1.167,2.55) -- (1.167,2.28);
\draw[polcarico] (1.333,2.55) -- (1.333,2.28);
\draw[polcarico] (1.500,2.55) -- (1.500,2.28);
\node[anchor=east,inner sep=3pt] at (0.5,2.415) {$p$};
\draw[polcostruzione] (1.5,2.28) -- (1.500,1.000);
\fill (1.000,1.000) circle (1.1pt);
\draw[sgrrot] ([shift={(200:0.24)}]1.0,1.0) arc (200:-20:0.24);
\node[anchor=north,inner sep=3pt] at (1.000,0.950) {$p\ell^2$};
\draw[polcostruzione] (1.0,2.55) -- (1.0,3.35);
\draw[polforza,black!45] (1.0,3.3) -- (1.0,2.72) node[pos=0.25,anchor=east,inner sep=3pt,black] {$p\ell$};
\draw[polrisultante] (2.0,3.3) -- (2.0,2.72) node[pos=0.25,anchor=west,inner sep=3pt,black] {$p\ell$};
\draw[polquota] (1.0,3.55) -- (2.0,3.55) node[midway,anchor=south,inner sep=2.5pt] {$\ell$};
\pic at (2.000,4.000) {polpunto};
\node[anchor=west,inner sep=4pt] at (2.000,4.000) {$K$};
\node at (1.25,-1.05) {(a)};
\end{scope}
\begin{scope}[shift={(6.3,4.9)},x=3.6cm,y=3.6cm]
\draw[polforza] (0.000,0.000) -- (0.000,-1.000) node[pos=0.5,anchor=west,inner sep=3pt] {$p\ell$};
\draw[polestremo] (0.000,-1.000) -- (-0.100,-0.200) node[pos=0.45,anchor=east,inner sep=3pt] {$\rb_B$};
\draw[polestremo] (-0.100,-0.200) -- (0.000,0.000) node[pos=0.5,anchor=east,inner sep=3pt] {$\rb_A$};
\pic at (0.000,0.000) {polpunto};
\pic at (0.000,-1.000) {polpunto};
\draw[fill=white,line width=.6pt] (-0.100,-0.200) circle (1.5pt);
\node[anchor=south,inner sep=3.5pt] at (0.000,0.000) {$0$};
\node[anchor=north,inner sep=3.5pt] at (0.000,-1.000) {$1$};
\node[anchor=north east,inner sep=1.5pt] at (-0.100,-0.200) {$m$};
\node at (0.0,-1.755) {(b)};
\end{scope}
\end{tikzpicture}%
